\section*{ID9405: Definition of Htri}
\paragraph{Metadata.}
PiID: 0906979396122 | RTEID: RTE:P36521.7036:T3.6279 | UUID: a421d96b-f80d-586e-a0f8-b9bb6cc860e2

\paragraph{Canonical Statement.}
cate Lean Certificate Coq Certificate Isabelle audit: included in aggregate HOL build. Lean audit: compiled corpus certificate. Rocq audit: compiled corpus certificate. Dependencies and Test Handle Dependencies [ID 0332], [ID 0300] Node registry\_id\_0333 Support Titles Variation Law; Support Relation; Closure Relation Test Handle If the cited entry equations cannot be made mutually admissible in one TRAWIN-normalized frame, the entry fails. Canonical Equation Links I3=𝐼 (𝐴:𝐵)+𝐼 (𝐴:𝐶 )−𝐼 (𝐴:𝐵𝐶 ) \$Htri=H𝐴⊗H𝐵⊗H𝐶\$ 𝜕I3=0 [ID 0334] Universal Avalanche Exponent TZPID ID0334 UUID 57a73aed-4ef1-5a9f-8d96-b60972c6f36b Type Transfer Statement Scale transfer LOC Classification Working routing: QA641 manifolds and topology; QC174.17 mathematical physics arXiv Classification Primary: math-ph; Cross-list: gr-qc, math.DG Canonical Statement This entry develops universal avalanche exponent inside the field theory and a

\paragraph{Canonical Equation.}
\[
Htri=H𝐴⊗H𝐵⊗H𝐶
\]

\paragraph{Dictionary.}
Definition of Htri — Htri=H𝐴⊗H𝐵⊗H𝐶

\paragraph{Encyclopedia.}
Definition of Htri is an algebraic definition, written Htri=H𝐴⊗H𝐵⊗H𝐶. It is an algebraic definition expressing the quantity directly in terms of the variables on its right-hand side. It is built from H, defining Htri. Within the Trawin Zero Point registry it serves as one element of the framework's mathematical narrative.
